\chapter{问题定义}

本章明确自进化强化学习系统的设计目标与核心问题，为后续方法设计奠定基础。

\section{设计目标}

本文的设计目标是构建一套面向智能体的自进化强化学习系统，使其在无人工标注条件下满足以下三个目标：

\textbf{目标一：数据自主产生。} 系统应能自主产生训练数据，而非依赖真实用户回流或人工采集。除初始种子外，后续所有训练任务由系统基于上一轮执行经验自主生成。

\textbf{目标二：奖励可信。} 奖励应锚定在环境真实状态而非智能体自述，从取证机制上抑制奖励作弊，使奖励信号可信。

\textbf{目标三：自进化闭环。} 系统应形成"采样—奖励—更新—生成下一轮任务"的闭环，持续自主运转，无需人工干预。

\section{核心问题}

围绕上述目标，本文需要解决以下核心问题：

\textbf{问题一：如何让系统自主产生下一轮训练任务？} 传统强化学习依赖固定数据集，训练任务由外部供给。自进化要求系统基于上一轮经验自主生成新任务，这涉及执行经验的概括、沙箱状态的跨步继承、以及新种子的生成。

\textbf{问题二：如何使奖励锚定在环境真实状态？} 若奖励仅看智能体自述的轨迹文本，模型会学会谎报完成。需要让奖励锚定在环境客观变更上，使作弊从机制上被抑制。

\textbf{问题三：如何保证自进化过程的稳定性？} 自进化系统持续自主运转，无人值守。策略一旦发生模式坍缩或训练崩溃而不被及时检测，将导致后续所有训练步在无效信号上空转。需要实时监控关键指标并在异常时自动终止。

\section{系统边界}

本文的系统边界如下：

\begin{itemize}
    \item 基座模型为 Qwen3.5-9B，策略优化采用标准 GRPO，不追求算法最优，仅为证明系统可行性。
    \item 沙箱环境为代码沙箱（本地或云端），支持文件系统操作与命令执行。
    \item 奖励裁判为外部冻结的大模型，以服务形式提供。
    \item 训练步数不设固定值，由终止判据决定。
\end{itemize}

\section{本章小结}

本章明确了系统的三个设计目标（数据自主、奖励可信、自进化闭环）与三个核心问题（自主产生任务、奖励锚定真实状态、保证稳定性），划定了系统边界。后续章节围绕这些问题展开方法设计与系统实现。
